\Askhsh{32390_SOLUTION}{1}
\begin{ylh}{1}
    \item [Ύλη από εικόνα]
\end{ylh}
ΛΥΣΗ
\begin{alist}
    \item Τα κρίσιμα σημεία της συνάρτησης $f$ στο διάστημα $[0,2]$, είναι τα εσωτερικά σημεία του διαστήματος στα οποία η $f$ δεν παραγωγίζεται ή η παράγωγός της είναι ίση με το μηδέν.
    Η συνάρτηση $f$ είναι παραγωγίσιμη άρα και συνεχής, με $f'(x) = 4x^3 - 4$ για κάθε $x \in [0,2]$. Είναι:
    \[ f'(x) = 0 \Leftrightarrow 4x^3 - 4 = 0 \Leftrightarrow x^3 = 1 \Leftrightarrow x = 1 \]
    Επομένως, η $f$ έχει ένα μόνο κρίσιμο σημείο, το $x_0 = 1$.
    \item Γνωρίζουμε ότι αν μία συνάρτηση $f$ είναι συνεχής σε ένα κλειστό διάστημα, τότε σύμφωνα με το θεώρημα μέγιστης και ελάχιστης τιμής, παρουσιάζει μέγιστο και ελάχιστο.
    Για την εύρεση του μέγιστου και του ελάχιστου εργαζόμαστε ως εξής:
    \begin{itemize}
        \item Βρίσκουμε τα κρίσιμα σημεία της $f$.
        \item Υπολογίζουμε τις τιμές της $f$ στα σημεία αυτά και στα άκρα του διαστήματος.
        \item Από αυτές τις τιμές, η μεγαλύτερη είναι το μέγιστο και η μικρότερη το ελάχιστο της $f$.
    \end{itemize}
    Οι τιμές της $f$ στο κρίσιμο σημείο της $x_0 = 1$ και στα άκρα του διαστήματος $[0,2]$ είναι
    \[ f(1) = -1, f(0) = 2 \text{ και } f(2) = 10. \]
    Άρα, η μέγιστη τιμή της $f$ στο $[0,2]$ είναι ίση με $10$ και η ελάχιστη τιμή της $f$ στο $[0,2]$ είναι ίση με $-1$.
\end{alist}